%%%% proposta.tex
%%%% Project specification paper based on the IJCAI--ECAI 26 template.

\typeout{Project paper using IJCAI--ECAI 26 formatting}

\documentclass{article}
\pdfpagewidth=8.5in
\pdfpageheight=11in

% Official IJCAI--ECAI 26 style file.
% Compile from this directory so the relative path resolves correctly.
\usepackage{FormattingGuidelines-IJCAI-ECAI-26/ijcai26}

% Packages required/recommended by the official template.
\usepackage{times}
\usepackage{soul}
\usepackage{url}
\usepackage[hidelinks]{hyperref}
\usepackage[utf8]{inputenc}
\usepackage[small]{caption}
\usepackage{graphicx}
\usepackage{amsmath}
\usepackage{amsthm}
\usepackage{booktabs}
\usepackage{algorithm}
\usepackage{algorithmic}
\usepackage[switch]{lineno}

% Keep line numbers enabled for a submission. Disable for camera-ready.
\linenumbers

\urlstyle{same}

% Optional theorem environments.
\newtheorem{example}{Example}
\newtheorem{theorem}{Theorem}

% Required PDF metadata block. Keep unchanged.
\pdfinfo{
/TemplateVersion (IJCAI.2026.0)
}

\title{Semantic segmentation of AIS trajectories for fishing activity detection using overlapping windows and temporal post-processing}

% For an anonymous submission, replace this block with the track-approved
% anonymous author declaration. For the camera-ready version, provide all
% author and affiliation information.
\author{
    Dinis Soveral Botelho Rosa
    \affiliations
    University of Coimbra / Department of Informatics Engineering
    \emails
    dinisrosa22@gmail.com
}

\begin{document}

\maketitle

% Project proposal: maximum of three pages in the IJCAI two-column format.
% The sections below correspond to the requirements in the course brief.

\section{Problem and Motivation}
\label{sec:problem-motivation}

Describe the problem addressed by the project, its context, and the people or
organizations affected by it. Explain why the problem is interesting and
important. Make explicit how the project relates to Human-Centered Artificial
Intelligence, including the relevant course topic or topics.

\section{Background}
\label{sec:background}

Review the scientific literature that provides context for the problem and
discuss existing approaches. Explain the limitations or open questions that
motivate the proposed project. Support the discussion with appropriate
bibliographic references, for example, \verb|\cite{key}|.

\section{Goals}
\label{sec:goals}

State the main aim of the project and the specific objectives that will be
pursued. Define what would count as a successful result and identify the
expected scientific, technical, or human-centered contributions.

\section{Approach}
\label{sec:approach}

Describe the planned methods or algorithms at a level appropriate for a
project proposal. Explain the intended system workflow and identify existing
implementations, libraries, or models that will be used, adapted, or
extended. State the main assumptions and design decisions.

\section{Materials and Data}
\label{sec:materials-data}

Describe the datasets, tools, frameworks, and other resources required by the
project. Explain how existing data will be prepared and, if new data is
needed, how it will be collected. Include relevant privacy, consent, fairness,
or other ethical considerations.

\section{Evaluation}
\label{sec:evaluation}

Explain how the proposed work will be evaluated. Identify the experimental
setup, baselines or comparison methods, and quantitative metrics or
qualitative analyses. State what outputs are expected, such as plots, tables,
user feedback, or case studies, and how they will support the conclusions.

% References may use the provided IJCAI named bibliography style.
\bibliographystyle{FormattingGuidelines-IJCAI-ECAI-26/named}
\bibliography{references}

\end{document}
